\section{Primitive forms and invariance within a genus}\label{sec:forms}

The principal class is not special for value counting. The orientation
choices used in Section~\ref{sec:norms} fill a compatible coset of
\(2\operatorname{Pic}(\mathcal O)\), rather than only the principal
class. We first make this precise for one discriminant.

Let \(\Delta=f^2D<0\), with \(D\) fundamental, and put
\(\mathcal O=\Z+f\mathcal O_K\), \(K=\Q(\sqrt D)\),
\(G=\operatorname{Pic}(\mathcal O)\). All forms in the first two
theorems are primitive positive definite integral binary quadratic
forms of discriminant \(\Delta\). For such a form \(Q\), define
\[
 \begin{aligned}
 P_Q&=\{Q(x,y):x,y\in\Z,\ \gcd(x,y)=1\},\\
 \mathscr V_{p,Q}&=Q(\Z_p^2\setminus p\Z_p^2),\qquad
 \Lambda_Q=\{n\ge1:n\in\mathscr V_{p,Q}\text{ for every }p\}.
 \end{aligned}
\]
Only positive values are included. Coordinates may vanish, and
each numerical value is counted once. Write
\[
 R_{T_1,T_2}(N)=\#\{(a,b)\in T_1\times T_2:a+b=N\},\qquad
 R_{Q_1,Q_2}=R_{P_{Q_1},P_{Q_2}}.
\]
Two forms are in the same genus when they are equivalent over every
\(\Z_p\) (and over \(\mathbb R\)). Such equivalences preserve primitive
vectors, hence their primitive local value sets.

Put \(S=\{p:p\mid2\Delta\}\), \(\chi=\chi_D\), and
\[
 J_\Delta(N)=\prod_{\substack{p\mid N\\p\notin S,\ \chi(p)=-1}}
                            \frac{p-1}{p-2}.
\]
For a pair \(Q_1,Q_2\), a target is locally admissible if
\(N\in\mathscr V_{p,Q_1}+\mathscr V_{p,Q_2}\) for every \(p\).
Set \(\eta_*(G)=+\infty\) when \(2G=0\); otherwise use the
prime-divisor threshold in Proposition~\ref{prop:norm-class-error}.

\begin{theorem}[Classical genus invariance]\label{thm:genus-binary-invariance}
Suppose \(Q_i,Q'_i\) are in the same genus for each \(i=1,2\).
For every fixed \(0<\eta<\eta_*(G)\), uniformly on their common
sufficiently large locally admissible targets,
\[
 R_{Q_1,Q_2}(N)
   =R_{Q'_1,Q'_2}(N)\left(1+O((\log N)^{-\eta})\right).
\]
Both counts have order \(NJ_\Delta(N)/\log N\). Nonadmissible
targets have zero count. If \(2G=0\), the two counts are equal
for every target. The constants and threshold depend on the fixed
forms and on \(\eta\). This theorem uses the classical route of
Appendix~\ref{sec:older-route} and the class-error sieve; it does not
use the recent zero-free or weighted-dilation inputs.
\end{theorem}

To state the explicit leading term, define the mixed local value density
\[
 d_{p;1,2}(N)=\int_{\Z_p}
    \1_{\mathscr V_{p,Q_1}}(a)\1_{\mathscr V_{p,Q_2}}(N-a)\,da
       \quad(p\in S),
\]
and the positive constant
\begin{equation}\label{eq:form-universal-factor}
 \mathfrak b_{D,S}=\frac{4L(1,\chi)}\pi
  \prod_{p\in S}\frac{1-\chi(p)/p}{1-1/p}
  \prod_{\substack{p\notin S\\\chi(p)=-1}}
                      \left(1-\frac2{p(p-1)}\right).
\end{equation}

\begin{theorem}[Mixed primitive form values]\label{thm:form-binary-leading}
Under the analytic inputs of Section~\ref{sec:inputs}, for fixed
\(Q_1,Q_2\) as above, uniformly through sufficiently large locally
admissible targets,
\begin{equation}\label{eq:form-binary-leading}
 R_{Q_1,Q_2}(N)\sim
   \mathfrak b_{D,S}\left(\prod_{p\in S}d_{p;1,2}(N)\right)
                  J_\Delta(N)\frac N{\log N}.
\end{equation}
Each bad-prime factor is a computable rational periodic function of
the target. The leading coefficient depends only on the two genera.
There is no class-number divisor in this formula.
\end{theorem}

The absence of such a divisor concerns numerical values, not weighted
representations by ideals or coordinates. A saturated value can be
represented by every class in each compatible genus. Selecting one
class therefore does not take a fixed fraction of a disjoint partition.

\subsection{The target ideal class and its local shadow}

For \(Q(x,y)=ax^2+bxy+cy^2\), take
\[
 \beta=(b+\sqrt\Delta)/2,\qquad
 A=\Z a+\Z\beta\subset\mathcal O,\qquad c_Q=-[A]\in G.
\]
Here \(\mathcal O=\Z[\beta]\), \(n(A)=a\), and
\(N(ax+\beta y)/a=Q(x,y)\). The ideal is proper invertible:
\[
 A\overline A=a(a,\beta,\overline\beta,c)\mathcal O=a\mathcal O,
\]
since \(\beta+\overline\beta=b\) and \(\gcd(a,b,c)=1\).
This construction allows \(a\) to share primes with the conductor.
It also fixes the class sign in our convention. For the classical
normalized ideal-form description compare
\cite[Lemma 10.2]{khayutin}.

\begin{lemma}\label{lem:form-ideal-shadow}
For every positive integer \(n\),
\begin{align*}
 n\in P_Q
 &\ \Longleftrightarrow\
 \text{some primitive integral invertible ideal }I
       \text{ has }n(I)=n,\ [I]=c_Q,\\
 n\in\Lambda_Q
 &\ \Longleftrightarrow\
 \text{some primitive integral invertible ideal }I
       \text{ has }n(I)=n,\ [I]\in c_Q+2G.
\end{align*}
\end{lemma}
\begin{proof}
An element \(\alpha\in A\) gives \(I=(\alpha)A^{-1}\), with
\(n(I)=N(\alpha)/a\) and class \(-[A]\). It is integral because
\((\alpha)\subset A\). Multiplication by the invertible ideal \(A\)
shows \(I\subset p\mathcal O\) if and only if \(\alpha\in pA\).
This is exactly the coordinate-primitivity condition. Conversely,
if \([I]=-[A]\), then \(IA=(\alpha)\subset A\), and the same
identities give a primitive representation of \(n(I)\).

Write \(A_p=z_p\mathcal O_p\). The norm identities in
\eqref{eq:norm-ideal-identities} give the units
\[
 \lambda_p=N(z_p)/a\in\Z_p^\times,\qquad
 \mathscr V_{p,Q}=\lambda_p\mathscr V_p,
\]
where \(\mathscr V_p\) is the principal primitive norm-value set for
\(\mathcal O_p\). The change from an integral basis of \(A\) to
the local generator is in \(\operatorname{GL}_2(\Z_p)\).
Moreover \(\Theta_p(c_Q)=\lambda_p\pmod{U_p}\), with \(\Theta\)
and \(U_p\) from Lemma~\ref{lem:norm-genus}.

If \(n\in\Lambda_Q\), choose primitive local elements
\(\alpha_p=z_pw_p\in A_p\) with \(N(\alpha_p)=an\).
The ideals \(w_p\mathcal O_p\) are primitive integral invertible,
and equal \(\mathcal O_p\) outside finitely many primes. Glue them
as in Lemma~\ref{lem:norm-genus}. Their norm valuations are
\(v_p(n)\), so the resulting ideal \(I\) has norm \(n\), and
\[
 \Theta_p([I])=n/N(w_p)=\lambda_p\pmod{U_p}.
\]
The kernel of \(\Theta\) is \(2G\), proving the coset assertion.
Conversely, for such an \(I\), write \(I_p=i_p\mathcal O_p\).
Equality of the \(\Theta\)-classes supplies a unit \(u_p\) with
\(N(u_p)=n/(\lambda_pN(i_p))\). Then
\(\alpha_p=z_pi_pu_p\) is primitive in \(A_p\) and has norm
\(an\). Outside \(S\) the order-unit norm group is all of
\(\Z_p^\times\), so the same conclusion holds there.
This proves the local-shadow equivalence at conductor primes as well.
\end{proof}

The coset condition is existential. A conductor-divisible value need
not determine a unique genus. For example, at discriminant \(-32\),
\(x^2+8y^2\) and \(3x^2+2xy+3y^2\) have primitive odd values
\(1\pmod8\) and \(3\pmod8\), respectively, and hence different
genera. Both nevertheless primitively represent \(24\), at
\((4,1)\) and \((3,-1)\).

\subsection{Saturation at an arbitrary class}

The saturation argument now targets \(c_Q\), rather than zero.
Start with the ideal in \(c_Q+2G\) from
Lemma~\ref{lem:form-ideal-shadow}. Initial reorientation of the
reserved exponent-one primes keeps its class in that coset.
If their doubled classes generate \(2G\), the independent reversal
ranges cancel the difference from \(c_Q\). They leave every
conductor-local ideal unchanged. Thus a failure in
\(\Lambda_Q\setminus P_Q\) has the same cover
\eqref{eq:norm-deficiency}, with exactly the same \(\Omega_M,K_M\).

For two forms, charge a pair in
\(\Lambda_{Q_1}\times\Lambda_{Q_2}\setminus(P_{Q_1}\times P_{Q_2})\)
to an exceptional summand. Outside \(S\), both summands are units
at every inert prime. Mark only the selected summand's deficient
family. Its marked split primes are disjoint from these inert primes,
and remove the same \(p-1\) roots modulo \(p^2\). The calculation
\eqref{eq:norm-marked-sieve} therefore gives, for every
\(0<\eta<\eta_*(G)\),
\begin{equation}\label{eq:form-class-error}
 0\le R_{\Lambda_{Q_1},\Lambda_{Q_2}}(N)-R_{Q_1,Q_2}(N)
   \ll_{\Delta,\eta}\frac{NJ_\Delta(N)}{(\log N)^{1+\eta}}
                  \qquad(N\ge3).
\end{equation}
When \(2G=0\), this difference is identically zero. In particular,
the exponent \(\eta=1/5\) remains available for every fixed order.
No recent analytic input enters this comparison.

\subsection{The separate genus filters}

The local genus biases must still be retained. Set
\[
 L=8\prod_{\substack{p\in S\\p\text{ odd}}}p,\qquad
 \alpha=\varphi(L)/L,\quad H=\ker(\chi\bmod L),\quad
 \mathcal C=\langle p:p\notin S,\ \chi(p)=1\rangle.
\]
For each \(S\)-supported \(r\), put
\[
 A_{r,i}=\{u\in(\Z/L\Z)^\times:
                      ru\in\mathscr V_{p,Q_i}\text{ for all }p\in S\}.
\]
These sets are well defined and are unions of unit-square states:
equal units modulo an odd \(p\), or modulo eight at two, differ
by a local unit square; multiplication of a primitive vector by a
rational unit scales its value by that square.

Crucially \(A_{r,i}\subset H\), although individual local genus
characters can be nonprincipal. To verify this, use an ideal \(A_i\)
as above, with norm \(a_i\) and local multipliers \(\lambda_{i,p}\).
At every prime outside \(S\), an ideal norm has even valuation at
an inert prime and has no field-norm obstruction at a split prime.
Positivity and Hilbert reciprocity therefore give
\[
 \prod_{p\in S}(\lambda_{i,p},D)_p
      =\prod_{p\in S}(a_i,D)_p=1.
\]
Here the numerator in \(\lambda_{i,p}=N(z_{i,p})/a_i\) is a
field norm. If \(u\in A_{r,i}\), then
\((ru,D)_p=(\lambda_{i,p},D)_p\) at every bad prime.
The positive \(S\)-supported number \(r\) has product of bad
symbols one. For an integer unit representative \(u\), reciprocity
identifies \(\prod_{p\in S}(u,D)_p\) with \(\chi(u)\).
Multiplication of the displayed local relations gives \(\chi(u)=1\).
This does not identify \(A_{r,i}\) with the principal-form filter.

Outside \(S\), the primitive value set is \(\Z_p\) at split primes
and \(\Z_p^\times\) at inert primes. Consequently the exact disjoint
decomposition is
\begin{equation}\label{eq:form-prefix-decomposition}
 \Lambda_{Q_i}=\coprod_{r\in\langle S\rangle}
                  r\{x\in\mathcal C:x\bmod L\in A_{r,i}\}.
\end{equation}
For \(Q_i\), the identity
\(\operatorname{adj}(\operatorname{Hess}Q_i)\nabla Q_i(v)=-\Delta v\)
shows that some partial derivative at a primitive vector has valuation
at most \(v_p(\Delta)\). Hensel lifting therefore makes
\(\mathscr V_{p,Q_i}\) a union of full cosets modulo
\(p^{e_p}\), \(e_p=2v_p(\Delta)+1\). Thus each \(d_{p;1,2}\)
is a rational periodic function, obtained by counting values in this
finite quotient. Its positive values have a fixed positive minimum.
At good primes, which are odd, every target is a sum of two allowed
local values. Local admissibility is precisely
\(\prod_{p\in S}d_{p;1,2}(N)>0\).

\subsection{Mixed profiles and their normalization}

For fixed prefixes \(r,s\), let \(g=(r,s)\). When \(g\mid N\), put
\(A=r/g\), \(B=s/g\), \(M=N/g\), and define
\[
 \Gamma_{r,s;1,2}(N)=\frac L{|H|^2}
  \#\{(u,v)\in A_{r,1}\times A_{s,2}:Au+Bv=M\bmod L\};
\]
otherwise set it to zero. The original two filters are retained after
dividing by \(g\).

There are finitely many profiles that cover every locally admissible
target. Indeed, for local summands \(a+b=N\), replace them by
\(a+p^{e_p}t,b-p^{e_p}t\). Values remain in their full local
cosets. Valuations below \(e_p\) stay fixed; otherwise each summand
excludes at most one residue \(t\bmod p^2\) if its valuation is to
be at most \(e_p+1\). Since \(p^2>2\), both bounds can be met.
Thus one can take
\(r,s\mid Z=\prod_{p\in S}p^{e_p+1}\).
CRT combines the local unit states into a root in
\(A_{r,1}\times A_{s,2}\). Conversely, a positive root lifts on the
local line, because one of \(A,B\) is a unit; the unit-square
conditions preserve each separate local value set.

Lemma~\ref{lem:older-quadratic-filter} allows distinct filters
\(U,V_0\subset H\), so it gives a positive lower bound of order
\(NJ_\Delta(N)/\log N\) in any positive profile. There are only
finitely many such profiles and root patterns. The unmarked inert
sieve supplies the matching upper bound for the whole mixed shadow.
Together with \eqref{eq:form-class-error}, this proves the classical
order bounds for \(R_{Q_1,Q_2}\).
If \(Q_i,Q'_i\) are in the same genus, their local value sets and
their mixed shadow coincide exactly. Subtract the two estimates
\eqref{eq:form-class-error}, then divide by the classical positive
lower bound. This proves Theorem~\ref{thm:genus-binary-invariance}.
If \(2G=0\), equality of the individual global and local sets gives
the exact conclusion for every target.

For the explicit leading term, apply Theorem~\ref{thm:general} to
each fixed profile. Write \(R_{r,s;1,2}(N)\) for the contribution
of prefixes \(r,s\) in \eqref{eq:form-prefix-decomposition}.
Let \(c_{\mathcal C}\) be its marginal core
constant and \(C_{\rm in}\) the good inert product in
Section~\ref{sec:norms}. The normalized contribution is
\[
 c_{\mathcal C}^2 C_{\rm in}\Gamma_{r,s;1,2}(N)\frac g{rs}+o(1).
\]
At a good inert prime the coefficients of the line \(rx+sy=N\)
are units, and its two forbidden roots merge exactly when \(p\mid N\).
Discarding the two residue filters in the upper sieve gives
\[
 R_{r,s;1,2}(N)\ll_\Delta
       \frac{NgJ_\Delta(N)}{rs\log N}+O(N^{9/20}).
\]
Only polynomially many logarithmic prefix pairs can occur, and
\(\sum_{r,s\in\langle S\rangle}(r,s)/(rs)<\infty\).
The finite-prefix limit and the uniform summable tail thus justify
summing the fixed-profile main terms.

For completeness, the mixed Haar normalization is
\begin{equation}\label{eq:form-haar-strata}
 \sum_{r,s}\Gamma_{r,s;1,2}(N)\frac g{rs}
        =\frac4{\alpha^2}\prod_{p\in S}d_{p;1,2}(N).
\end{equation}
At a prime with \(v_p(r)=i\), \(v_p(s)=j\), and \(v_p(L)=k\),
each compatible unit-pair root modulo \(p^k\) has measure
\(p^{-\max(i,j)-k}\) on the line. If \(i\le j\), prescribing the
second summand modulo \(p^{j+k}\) fixes the first summand's required
residue; the other case is symmetric. Hence a global root has
measure \(g/(Lrs)\). Since \(|H|=\varphi(L)/2\), this is
\((\alpha^2/4)\Gamma_{r,s;1,2}(N)g/(rs)\) after summing its roots.
Summing the disjoint valuation strata proves
\eqref{eq:form-haar-strata}. The Euler identity
\eqref{eq:norm-core-constant} and the good-prime cancellation in
Section~\ref{sec:norms} now give \eqref{eq:form-universal-factor}.
Finally \eqref{eq:form-class-error} transfers the shadow asymptotic
to the actual values, proving Theorem~\ref{thm:form-binary-leading}.
The class group has entered only this smaller error, not a divisor
of the leading value count.

\subsection{Different fixed conductors in one field}

\begin{corollary}\label{cor:common-field-form-pairs}
Fix an imaginary quadratic field of fundamental discriminant \(D\),
two orders of discriminants \(\Delta_i=f_i^2D\), and primitive positive
definite forms \(Q_i\) of discriminants \(\Delta_i\). Put
\[
 S=\{p:p\mid2\Delta_1\Delta_2\},\qquad
 J_S(N)=\prod_{\substack{p\mid N\\p\notin S,\ \chi_D(p)=-1}}
                       \frac{p-1}{p-2}.
\]
The mixed classical order bounds hold on sufficiently large locally
admissible targets. Under the inputs of Section~\ref{sec:inputs},
\[
 R_{Q_1,Q_2}(N)\sim
 \mathfrak b_{D,S}\left(\prod_{p\in S}d_{p;1,2}(N)\right)
                         J_S(N)\frac N{\log N}.
\]
If each \(Q'_i\) is in the genus of \(Q_i\), then, independently
of those recent inputs, the classical comparison is
\[
 R_{Q_1,Q_2}(N)=R_{Q'_1,Q'_2}(N)
                         \left(1+O((\log N)^{-\eta})\right)
\]
for every
\(0<\eta<\min_i\eta_*(\operatorname{Pic}(\mathcal O_i))\), on
sufficiently large common locally admissible targets. If both doubled
class groups vanish the counts agree exactly; \(\eta=1/5\) is
otherwise a common permissible exponent. All data remain fixed.
\end{corollary}
\begin{proof}
Use the common good split-prime core outside \(S\), the same quadratic
character \(\chi_D\), and \(L=8\prod_{p\in S,\ p\text{ odd}}p\).
Put \(\alpha=\varphi(L)/L\) and \(H=\ker(\chi_D\bmod L)\).
For each ideal lattice the local multipliers remain units, and their
Hilbert symbols have product one over \(S\); hence the two separate
prefix filters still lie in \(H\). Use
\(e_p=1+2\max_i v_p(\Delta_i)\) in the local finite quotient and
bounded-valuation profile construction. The root formula and
\eqref{eq:form-haar-strata} then use these two filters and the same
factor \(4/\alpha^2\).

Apply orientation saturation inside each order's own Picard group.
Deleting \(S\) from its prime families removes only finitely many
primes, so their Mertens densities are unchanged. For a bad pair mark
one summand at a time, with that order's maximal-subgroup family.
Both summands still exclude the same good inert primes. Thus the
same binary sieve gives the two class errors at their respective
thresholds, and their sum at the minimum threshold. The two Picard
groups need never be compared or multiplied, and overlapping prime
families across the orders cause no problem: no sieve marks both
partitions simultaneously. No joint ring-class prime-distribution
theorem is required.

The classical bounds and genus comparison use
Lemma~\ref{lem:older-quadratic-filter}, while the leading asymptotic
uses Theorem~\ref{thm:general} on each fixed profile. The uniform
prefix majorant and summable Haar strata use only the common core
and separate masks.
Its Euler identity is \eqref{eq:norm-core-constant} with the enlarged
\(S\), giving precisely \(\mathfrak b_{D,S}\). This proves the stated
asymptotic and classical comparison. Genus replacement is made in
each respective discriminant; no cross-discriminant identification
of genera or assertion for different quadratic fields is involved.
\end{proof}
